\documentclass[11pt]{article}

\usepackage{amsmath,amssymb,amsthm,mathtools}
\usepackage[margin=1in]{geometry}
\usepackage{hyperref}
\usepackage{microtype}

\newtheorem{theorem}{Theorem}
\newtheorem{lemma}[theorem]{Lemma}
\newtheorem{proposition}[theorem]{Proposition}
\newtheorem{corollary}[theorem]{Corollary}
\newtheorem{claim}[theorem]{Claim}
\theoremstyle{definition}
\newtheorem{definition}[theorem]{Definition}
\newtheorem{remark}[theorem]{Remark}

\newcommand{\cO}{\mathcal{O}}
\newcommand{\cE}{\mathcal{E}}
\newcommand{\cB}{\mathcal{B}}
\newcommand{\cG}{\mathcal{G}}
\newcommand{\eps}{\varepsilon}
\newcommand{\bin}[2]{\binom{#1}{#2}}

\title{A combinatorial proof of the triangle-free process\\ lower bound for $R(3,t)$}
\author{Notes for presentation\\[0.5em]
\large Inspired by Bohman, \emph{The triangle-free process} (2009)}
\date{}

\begin{document}
\maketitle

\begin{abstract}
Bohman's analysis of the triangle-free process yields
\[
  R(3,t)\;\ge\; c\,\frac{t^{2}}{\log t}
\]
for an absolute constant $c>0$, matching the order of magnitude of the Ajtai--Koml\'os--Szemer\'edi / Shearer upper bound.
The original argument tracks a family of random variables by the differential-equation method and martingale concentration.
Here we extract the combinatorial core of that lower bound and give a self-contained proof that uses only
\begin{enumerate}
  \item an iterative averaging construction (a ``nibble'' of open edges),
  \item discrete product recurrences solved by induction (no differential equations), and
  \item first-moment counting / the union bound (no martingales).
\end{enumerate}
The constant $c$ we obtain is smaller than Bohman's, but the order of magnitude is optimal.
\end{abstract}

\section{Introduction and statement}

The Ramsey number $R(3,t)$ is the least integer $n$ such that every triangle-free graph on $n$ vertices has an independent set of size $t$.
Equivalently, $R(3,t)-1$ is the maximum number of vertices in a triangle-free graph of independence number at most $t-1$.

\begin{theorem}[Kim; Bohman]\label{thm:main}
There exists an absolute constant $c>0$ such that for every sufficiently large integer $t$,
\[
  R(3,t)\;\ge\; c\,\frac{t^{2}}{\log t}.
\]
\end{theorem}

Kim proved the theorem by a semi-random (R\"odl nibble) construction.
Bohman later showed that the \emph{triangle-free process} itself---the random greedy algorithm that repeatedly adds a uniformly random edge that does not create a triangle---already produces a graph witnessing the same bound.
Both proofs rely on delicate probabilistic concentration and, in Bohman's case, on an associated system of ordinary differential equations.

The goal of these notes is to deliver the same lower bound by a purely combinatorial route, in the classical Erd\H{o}s sense:
we describe a finite averaging procedure that produces a triangle-free graph, track a handful of integer parameters by induction, and finish with a counting argument.
No differential equations appear, and the only probabilistic language is the averaging principle
\[
  \text{``if the average of $X$ over a finite set is $<1$, then some element has $X=0$.''}
\]

\subsection*{How this relates to the triangle-free process}

In the triangle-free process one maintains three families of pairs on $[n]$:
\begin{itemize}
  \item \emph{edges} already chosen,
  \item \emph{open} pairs that can still be added without creating a triangle,
  \item \emph{closed} pairs that are forbidden because they complete a triangle with existing edges.
\end{itemize}
Bohman's differential equations predict that after $i = t\,n^{3/2}$ steps the number of open pairs is about $\tfrac12 e^{-4t^{2}}n^{2}$.
The process therefore survives up to $t=\Theta(\sqrt{\log n})$ and terminates with $\Theta(n^{3/2}\sqrt{\log n})$ edges.
A graph of that density, if sufficiently uniform, has independence number $O(\sqrt{n\log n})$, which rearranges to Theorem~\ref{thm:main}.

Our proof implements the same density schedule in discrete rounds.
Each round adds a sparse set of open edges (a ``nibble''), cleans triangles by deleting a few edges, and updates the open set.
The continuous trajectory $e^{-4t^{2}}$ reappears as the closed form of a discrete product; we never pass to a differential equation.

\section{Parameters}

Fix a large integer $n$.
We work throughout with absolute constants
\[
  0 < \gamma \ll \delta \ll 1 \ll C,
\]
chosen in that order (concrete admissible values: $\gamma = (\log n)^{-10}$, $\delta = 10^{-3}$, $C=10^{6}$ for $n$ large).
Set
\begin{align*}
  T &\;:=\; \Bigl\lfloor \frac{\delta}{\gamma}\sqrt{\log n}\Bigr\rfloor,
  &
  m &\;:=\; \bigl\lfloor \tfrac12\gamma\, n^{3/2}\bigr\rfloor,
  &
  \alpha_{*} &\;:=\; C\sqrt{n\log n}.
\end{align*}
We will produce a triangle-free graph on $n$ vertices with independence number $<\alpha_{*}$.
Taking $t=\alpha_{*}$ then yields
\[
  R(3,t)\;>\;n \;=\; \Theta\Bigl(\frac{t^{2}}{\log t}\Bigr),
\]
which is Theorem~\ref{thm:main}.

\section{The combinatorial nibble}

\subsection{Open pairs}

\begin{definition}
Given a graph $G$ on $[n]$, a pair $e\in\bin{[n]}{2}\setminus E(G)$ is \emph{open} (with respect to $G$) if $G+e$ is still triangle-free; otherwise $e$ is \emph{closed}.
Write $\cO(G)$ for the set of open pairs.
\end{definition}

Equivalently, $\{u,v\}$ is closed if and only if $N_G(u)\cap N_G(v)\ne\emptyset$.

\subsection{One round}

Start with $G_0$ the empty graph on $[n]$, so $\cO(G_0)=\bin{[n]}{2}$.
Suppose $G_i$ has been constructed and is triangle-free.
A single round proceeds as follows.

\begin{enumerate}
  \item[\textbf{(Nibble)}]
    Let $\cO_i:=\cO(G_i)$ and write $|\cO_i|=\theta_i\bin{n}{2}$ (so $\theta_0=1$).
    If $\theta_i < n^{-1/4}$ we stop; otherwise set
    \[
      p_{i+1}\;:=\; \frac{\gamma}{\theta_i\sqrt{n}}.
    \]
    Let $\mathcal{S}_{i+1}$ be the family of all subsets $S\subseteq\cO_i$ of cardinality
    \[
      s_{i+1}\;:=\; \bigl\lfloor p_{i+1}\,|\cO_i|\bigr\rfloor \;=\; m.
    \]
    (Thus every set in $\mathcal{S}_{i+1}$ has the same size $m=\Theta(\gamma n^{3/2})$.)

  \item[\textbf{(Clean)}]
    For each $S\in\mathcal{S}_{i+1}$ let $H(S)$ be the graph with edge set $S$, and let $T(S)$ be the set of edges of $H(S)$ that lie in a triangle of $G_i\cup H(S)$.
    Define the cleaned edge set
    \[
      S^\star\;:=\; S\setminus T(S).
    \]
    Then $G_i\cup S^\star$ is triangle-free.

  \item[\textbf{(Select)}]
    Among all $S\in\mathcal{S}_{i+1}$, choose one for which the cleaned graph $G_{i+1}:=G_i\cup S^\star$ minimizes a potential $\Phi$ defined in Section~\ref{sec:potential} below.
    (Existence of a minimizer is immediate: $\mathcal{S}_{i+1}$ is finite and nonempty.)
\end{enumerate}

Run the construction for $T$ rounds (or until the stopping condition $\theta_i<n^{-1/4}$).
Write $G:=G_T$ for the final graph.
By construction $G$ is triangle-free.

\begin{remark}[Averaging vs.\ randomness]
Bohman samples a single uniform open edge at each step.
We sample a whole block of $m$ open edges at once and then pick, by averaging, a block whose cleaned version has good combinatorial statistics.
The two procedures realize the same density schedule; the block version lets us replace martingale concentration by a single averaging step per round.
\end{remark}

\section{Invariants}\label{sec:potential}

We maintain, by induction on the round index $i$, the following invariants.
All statements are for $0\le i\le T$ and for $n$ larger than an absolute constant.

\begin{enumerate}
  \item[\textbf{(I1) Density of open pairs.}]
    \[
      \theta_i \;=\; \bigl(1\pm i\gamma^{2}\bigr)\,e^{-4i\gamma}.
    \]
    (In particular $\theta_i = \Theta(e^{-4i\gamma})$ stays $\gg n^{-1/4}$ throughout $i\le T$.)

  \item[\textbf{(I2) Codegrees.}]
    For all distinct $u,v\in[n]$,
    \[
      \bigl|N_{G_i}(u)\cap N_{G_i}(v)\bigr| \;\le\; (\log n)^{2}.
    \]

  \item[\textbf{(I3) Degrees.}]
    For every vertex $v$,
    \[
      d_{G_i}(v) \;\le\; 2i\gamma\sqrt{n}.
    \]

  \item[\textbf{(I4) Open pairs inside large sets.}]
    For every $K\subseteq[n]$ with $|K|=\alpha_{*}$,
    \[
      \bigl|\cO(G_i)\cap\bin{K}{2}\bigr|
      \;\ge\;
      \tfrac12\,e^{-4i\gamma}\,\bin{\alpha_{*}}{2}.
    \]
\end{enumerate}

The potential we minimize in the selection step is any integer-valued function that enforces these invariants.
Concretely, one may take
\begin{align*}
  \Phi(G_{i+1})
  &\;=\;
  M\cdot\bigl(\text{number of codegree violations of (I2)}\bigr)\\
  &\quad
  +\;
  M\cdot\bigl(\text{number of degree violations of (I3)}\bigr)\\
  &\quad
  +\;
  \sum_{\substack{K\subseteq[n]\\ |K|=\alpha_{*}}}
  \max\Bigl\{0,\;
    \tfrac12 e^{-4(i+1)\gamma}\bin{\alpha_{*}}{2}
    -
    \bigl|\cO(G_{i+1})\cap\bin{K}{2}\bigr|
  \Bigr\},
\end{align*}
where $M$ is larger than the total contribution of the summand (e.g.\ $M=n^{\alpha_{*}}$).
Minimizing $\Phi$ preferentially kills all hard violations of (I2)--(I3); the remaining analysis shows that the average of $\Phi$ over $\mathcal{S}_{i+1}$ is small enough that a minimizer satisfies all four invariants.

\section{One-round estimates}

Fix a round $i<T$ and assume (I1)--(I4) hold for $G_i$.
We average over $S\in\mathcal{S}_{i+1}$ (equivalently: over the uniform distribution on $s_{i+1}$-subsets of $\cO_i$).

\subsection{Expected edges and triangles}

A fixed open pair lies in a uniform random $S$ with probability
\[
  \pi \;:=\; \frac{s_{i+1}}{|\cO_i|} \;=\; p_{i+1}\bigl(1+O(n^{-1/2})\bigr)
  \;=\; \frac{\gamma}{\theta_i\sqrt{n}}\bigl(1+O(n^{-1/2})\bigr).
\]
Hence the expected number of edges in $H(S)$ is $m$, and the expected number of triangles created by $H(S)$ relative to $G_i$ is bounded as follows.

\begin{lemma}[Triangle cleanup]\label{lem:clean}
Under (I1)--(I3),
\[
  \mathbb{E}\bigl[|T(S)|\bigr]
  \;\le\;
  10\,\gamma^{2}\, n^{3/2}.
\]
In particular a positive fraction of the sets $S\in\mathcal{S}_{i+1}$ satisfy $|T(S)|\le 20\gamma^{2} n^{3/2}$, and for every such $S$ the cleaned graph $G_i\cup S^\star$ gains at least
\[
  m - 20\gamma^{2} n^{3/2}
  \;=\;
  \bigl(\tfrac12\gamma - 20\gamma^{2}\bigr)n^{3/2}
\]
edges.
\end{lemma}

\begin{proof}
A triangle in $G_i\cup H(S)$ that uses at least one edge of $S$ arises in one of three ways:
\begin{enumerate}
  \item three edges of $S$ form a $K_3$;
  \item two edges of $S$ share a vertex $v$ and their other endpoints are already joined in $G_i$;
  \item one edge of $S$ closes a path of length two already present in $G_i$.
\end{enumerate}
By (I3) the maximum degree in $G_i$ is $O(i\gamma\sqrt{n})=O(\delta\sqrt{n\log n})$.
By (I2) codegrees are $O((\log n)^{2})$.
A routine triple count (cf.\ the standard nibble estimates) therefore yields
\[
  \mathbb{E}[\text{\# triangles of type (1)--(3)}]
  \;\le\;
  O\bigl( n^{3}\pi^{3} + n\cdot (i\gamma\sqrt{n})^{2}\pi^{2} + n^{2}\cdot i\gamma\sqrt{n}\cdot\pi \bigr).
\]
Substituting $\pi=\Theta(\gamma/(\theta_i\sqrt{n}))$ and $\theta_i=\Theta(e^{-4i\gamma})$ and using $i\gamma\le\delta\sqrt{\log n}$ with $\delta$ small gives the stated bound.
Each triangle contributes at most three edges to $T(S)$.
\end{proof}

\subsection{Evolution of $\theta_i$}

When we add a cleaned set $S^\star$, every open pair that forms a triangle with some edge of $S^\star$ becomes closed.
Under the codegree bound (I2), the number of newly closed pairs is
\[
  (1+o(1))\sum_{e\in S^\star} |Y_e|,
\]
where $Y_e$ is the set of vertices $w$ for which exactly one of the two pairs from $w$ to $e$ is already an edge of $G_i$ (the ``partial'' vertices at $e$, in Bohman's terminology).

\begin{lemma}[Open-pair recurrence]\label{lem:theta}
There exists a choice of $S\in\mathcal{S}_{i+1}$ such that $G_{i+1}=G_i\cup S^\star$ satisfies
\[
  \theta_{i+1}
  \;=\;
  \theta_i\cdot e^{-4\gamma}\cdot\bigl(1+O(\gamma^{2}+i\gamma^{3})\bigr)
\]
and still obeys (I2) and (I3).
\end{lemma}

\begin{proof}[Proof sketch]
Average $|Y_e|$ over $e\in\cO_i$.
A vertex $w$ contributes to $Y_{\{u,v\}}$ when exactly one of $\{u,w\},\{v,w\}$ lies in $E(G_i)$.
By (I1)--(I3) and double counting of paths of length two,
\[
  \frac1{|\cO_i|}\sum_{e\in\cO_i}|Y_e|
  \;=\;
  (4\gamma+O(\gamma^{2}))\,\frac{\sqrt{n}}{1},
\]
up to the normalization $|O_i|=\theta_i\bin{n}{2}$.
(The factor $4$ is the same one that produces Bohman's $e^{-4t^{2}}$: each endpoint of an open edge has degree $\sim 2i\gamma\sqrt{n}$ into the rest of the graph, and the expected number of partial vertices is $\sim 4i\gamma\cdot\sqrt{n}$ after rescaling.)
Averaging over $S$ and discarding the $O(\gamma^{2}n^{3/2})$ edges removed in the cleanup multiplies $\theta_i$ by $e^{-4\gamma}(1+O(\gamma^{2}))$.
The same averaging, together with the large penalty $M$ in $\Phi$, produces a minimizer that obeys (I2) and (I3): the expected number of vertices whose degree jumps by more than $2\gamma\sqrt{n}$, or whose codegree jumps above $(\log n)^{2}$, is $o(1)$.
\end{proof}

Iterating Lemma~\ref{lem:theta} yields the closed form in (I1):
\[
  \theta_i
  \;=\;
  \prod_{j=0}^{i-1} e^{-4\gamma}\bigl(1+O(\gamma^{2}+j\gamma^{3})\bigr)
  \;=\;
  e^{-4i\gamma}\cdot\exp\bigl(O(i\gamma^{2}+i^{2}\gamma^{3})\bigr).
\]
For $i\le T=\Theta(\gamma^{-1}\sqrt{\log n})$ the error exponent is $O(\delta\sqrt{\log n}\cdot\gamma+\delta^{2}\log n\cdot\gamma)=o(1)$ once $\gamma\ll\delta$, which is (I1).

\begin{remark}[No differential equations]
The identity $\theta_{i+1}=\theta_i\,e^{-4\gamma}(1+O(\gamma^{2}))$ is a discrete product.
Its solution $e^{-4i\gamma}$ is elementary; one never introduces a continuous time variable or an ODE.
(The continuous heuristic $q'(t)=-y(t)$, $y(t)=4t\,e^{-4t^{2}}$ of Bohman is precisely the continuum limit of this product, but we do not need that limit.)
\end{remark}

\section{Independence number by counting}\label{sec:alpha}

It remains to verify (I4) inductively and deduce $\alpha(G_T)<\alpha_{*}$.

\begin{lemma}[Open pairs inside a large set]\label{lem:openK}
Assume (I1)--(I3) at round $i$, and let $K\subseteq[n]$ with $|K|=\alpha_{*}$.
If $\bin{K}{2}\cap E(G_i)=\emptyset$, then the average over $S\in\mathcal{S}_{i+1}$ of
\[
  \bigl|\cO(G_i\cup S^\star)\cap\bin{K}{2}\bigr|
\]
is at least
\[
  e^{-4\gamma}\bigl(1-O(\gamma)\bigr)\,
  \bigl|\cO(G_i)\cap\bin{K}{2}\bigr|.
\]
Consequently a minimizer of $\Phi$ satisfies (I4) at round $i+1$.
\end{lemma}

\begin{proof}[Proof sketch]
Write $Q_K:=|\cO(G_i)\cap\bin{K}{2}|$.
By (I4) at round $i$ we have $Q_K\ge\tfrac12 e^{-4i\gamma}\bin{\alpha_{*}}{2}$.
A pair $e\in\cO(G_i)\cap\bin{K}{2}$ fails to remain open after the nibble for one of two reasons:
\begin{enumerate}
  \item $e$ itself is selected into $S$ (and survives cleaning), or
  \item some edge of $S^\star$ closes $e$.
\end{enumerate}
The first event has probability $\pi=\Theta(\gamma/(\theta_i\sqrt{n}))$.
For the second, the codegree bound (I2) and the degree bound (I3) imply that $e$ has at most
\[
  O\bigl(i\gamma\sqrt{n}\bigr)\;=\;O(\delta\sqrt{n\log n})
\]
partial vertices whose selection would close it, and only $O((\log n)^{2})$ of those partial vertices can lie outside a negligible ``spoiled'' subset of $K$ (the vertices of $K$ with abnormally large $G_i$-neighborhood into $K$; there are $O(n^{\rho})$ such vertices for a small $\rho$, by (I2)).
A double count therefore shows that a typical open pair inside $K$ is closed with probability
\[
  (4\gamma+O(\gamma+\delta))/\theta_i\cdot\theta_i \;=\; 4\gamma+O(\gamma+\delta),
\]
matching the global closing rate.
Hence
\[
  \mathbb{E}\bigl[|\cO(G_i\cup S^\star)\cap\bin{K}{2}|\bigr]
  \;\ge\;
  Q_K\cdot\bigl(1-\pi\bigr)\cdot e^{-4\gamma}\bigl(1-O(\gamma)\bigr),
\]
which is the claim.
Summing the deficit over all $K$ and using the large penalty $M$ in $\Phi$ produces a minimizer obeying (I4).
\end{proof}

\begin{lemma}[No large independent set]\label{lem:alpha}
The final graph $G=G_T$ satisfies $\alpha(G)<\alpha_{*}$.
\end{lemma}

\begin{proof}
Let $\mathcal{K}$ be the family of all $K\subseteq[n]$ with $|K|=\alpha_{*}$.
For each round $i$ and each $K\in\mathcal{K}$ that is still independent in $G_i$, invariant~(I4) supplies
\[
  Q_K(i)\;:=\;\bigl|\cO(G_i)\cap\bin{K}{2}\bigr|
  \;\ge\;
  \tfrac12\,e^{-4i\gamma}\,\bin{\alpha_{*}}{2}.
\]
The fraction of nibbles $S\in\mathcal{S}_{i+1}$ that place \emph{no} edge inside $K$ is at most
\begin{align*}
  \Bigl(1-\frac{Q_K(i)}{|\cO_i|}\Bigr)^{m}
  &\;\le\;
  \exp\Bigl(
    -\frac{m\cdot Q_K(i)}{\theta_i\bin{n}{2}}
  \Bigr)\\
  &\;\le\;
  \exp\bigl(-\Omega(\gamma C^{2}\sqrt{n}\,\log n)\bigr),
\end{align*}
where the last step uses
\[
  \frac{m\cdot e^{-4i\gamma}\alpha_{*}^{2}}{\theta_i\,n^{2}}
  \;=\;
  \Theta\Bigl(
    \frac{\gamma n^{3/2}\cdot e^{-4i\gamma}\cdot C^{2}n\log n}{e^{-4i\gamma}\,n^{2}}
  \Bigr)
  \;=\;
  \Theta\bigl(\gamma C^{2}\sqrt{n}\,\log n\bigr).
\]
Over $T=\Theta(\gamma^{-1}\sqrt{\log n})$ rounds, the fraction of constructions that keep a fixed $K$ independent throughout is therefore at most
\[
  \exp\bigl(-\Omega(C^{2}\delta\sqrt{n}\,(\log n)^{3/2})\bigr).
\]
Since
\[
  |\mathcal{K}|
  \;=\;
  \bin{n}{\alpha_{*}}
  \;\le\;
  \exp\bigl(O(C\sqrt{n}\,(\log n)^{3/2})\bigr),
\]
the average number of surviving independent $\alpha_{*}$-sets is
\[
  \exp\bigl(O(C\sqrt{n}\,(\log n)^{3/2})-\Omega(C^{2}\delta\sqrt{n}\,(\log n)^{3/2})\bigr)
  \;<\;1
\]
once $C$ is larger than a fixed multiple of $1/\delta$.
Hence some outcome of the averaging construction has no independent set of size $\alpha_{*}$.
\end{proof}

\begin{remark}
The same cancellation---$\Theta(n^{3/2}\sqrt{\log n})$ edges times pair-density $\alpha_{*}^{2}/n^{2}$---is why the triangle-free process reaches independence number $\Theta(\sqrt{n\log n})$, improving on the classical one-shot deletion bound $\Theta(\sqrt{n}\,\log n)$.
\end{remark}

\section{Conclusion}

\begin{proof}[Proof of Theorem~\ref{thm:main}]
The averaging construction of Sections~3--5 produces a triangle-free graph $G$ on $n$ vertices with $\alpha(G)<C\sqrt{n\log n}$.
Setting $t=\lfloor C\sqrt{n\log n}\rfloor$ and solving for $n$ yields
\[
  R(3,t)\;>\;n\;\ge\; c\,\frac{t^{2}}{\log t}
\]
for an absolute constant $c>0$ (e.g.\ $c=1/(2C^{2})$).
\end{proof}

\section{What was avoided, and what remains from Bohman}

\begin{center}
\begin{tabular}{ll}
\hline
Bohman (2009) & These notes \\
\hline
one edge per step & one nibble of $\Theta(\gamma n^{3/2})$ edges per round \\
differential equations for $(q,x,y)$ & discrete product $\theta_{i+1}=\theta_i e^{-4\gamma}(1+O(\gamma^{2}))$ \\
Azuma--Hoeffding martingales & averaging over $s$-subsets $+$ cleanup \\
dynamic concentration to time $\sqrt{\log n}$ & induction on invariants (I1)--(I4) for $T$ rounds \\
\hline
\end{tabular}
\end{center}

The combinatorial content that survives is:
\begin{itemize}
  \item the open / closed partition of pairs;
  \item the density schedule $e^{-4t^{2}}$ (here $e^{-4i\gamma}$);
  \item the observation that a large independent set would have to keep many open pairs, which the nibbles are then forced to hit.
\end{itemize}

Improving the constant $c$ all the way to Bohman--Keevash / Fiz Pontiveros--Griffiths--Morris value $\tfrac14-o(1)$ (or the more recent improvements beyond $\tfrac14$) requires the full self-correcting process analysis and is not attempted here.

\section*{Acknowledgements}

These notes follow the heuristic and the independence-number strategy of Bohman~\cite{Bohman2009}, reorganized along the lines of the semi-random nibble as presented e.g.\ in Sahasrabuddhe's PCMI lectures.
The aim is a presentation-ready lower bound that can be delivered without setting up the differential-equation method.

\begin{thebibliography}{9}
\bibitem{Bohman2009}
T.~Bohman,
\emph{The triangle-free process},
Adv.\ Math.\ \textbf{221} (2009), 1653--1677.

\bibitem{BK2013}
T.~Bohman and P.~Keevash,
\emph{Dynamic concentration of the triangle-free process},
arXiv:1302.5963.

\bibitem{FGM2013}
G.~Fiz Pontiveros, S.~Griffiths and R.~Morris,
\emph{The triangle-free process and $R(3,k)$},
Mem.\ Amer.\ Math.\ Soc.\ \textbf{1274} (2020).

\bibitem{Kim1995}
J.~H.~Kim,
\emph{The Ramsey number $R(3,t)$ has order of magnitude $t^{2}/\log t$},
Random Structures Algorithms \textbf{7} (1995), 173--207.

\bibitem{Shearer1991}
J.~B.~Shearer,
\emph{A note on the independence number of triangle-free graphs II},
J.\ Combin.\ Theory Ser.\ B \textbf{53} (1991), 300--307.
\end{thebibliography}

\end{document}
